\section{矿产资源法}
\begin{enumerate}
  \item 可考可不考。
  \item 矿产资源权属。国家所有权。
  \item 矿业权。
        \begin{enumerate}
          \item \zd{探矿权，采矿权分别}申请，登记。采矿 + 自己区域内为生产勘察不需要额外申请探矿证。
          \item 有偿取得，税费并行。资源税 \yjt 给国家，财政收入；资源补偿费 \yjt 给当地恢复环境。
          \item 转让。谁勘探，谁优先取得采矿权。完整勘探 + 最低投入。
        \end{enumerate}
  \item \zd{开采审批管理}。中央和省级有权审批。
        \begin{enumerate}
          \item 中央。特定矿种；海域；大型；矿区。
          \item 省。排除法。
        \end{enumerate}
  \item 不能开矿的地区。常识 + \zd{重要公路/河流/风景区}、一定距离以内不能移动的文物/名胜古迹。
  \item 集体矿山企业和个体采矿。
        \begin{enumerate}
          \item 集体 \yjt 鼓励和扶持。
          \item 个人 \yjt 严格限制。零星分散；普通建筑用的砂石、粘土；生活自用（煤）。
        \end{enumerate}
\end{enumerate}
